\section{What is known about satisfiability}
\label{sec:thresholds}

\subsection{Elementary satisfiability and first-moment bounds}
There is a sharp guarantee that applies to every proper $k$-CNF, whether
random or deterministic; see \citet[Theorem A, p.~4]{bangjensen2023}.
\begin{proposition}[A sharp universal clause-count guarantee]
\label{prop:fewclauses}
Let $n\ge k\ge2$ and $m\ge0$ be integers. If $F$ has $m$ clauses, each
containing exactly $k$ distinct variables chosen from $n$ available variables,
then its number $Z(F)$ of satisfying assignments obeys
\begin{equation}
 Z(F)\ge 2^n-m2^{n-k}=2^{n-k}(2^k-m).
 \label{eq:fewclauses}
\end{equation}
In particular, every such formula is SAT if $m<2^k$. With unused variables
and repeated clauses allowed, this condition is sharp: for every $m\ge2^k$
there is an UNSAT formula with these parameters.
\end{proposition}
\begin{proof}
A clause is false precisely when its $k$ literals are all false. This fixes
the values of its $k$ variables and leaves $n-k$ variables free, so it
excludes exactly $2^{n-k}$ assignments. The union of the sets excluded by
the $m$ clauses has size at most $m2^{n-k}$. Subtracting from $2^n$ proves
\eqref{eq:fewclauses}, whose right-hand side is positive when $m<2^k$.
No randomness or independence of the clauses is required.

For sharpness, take all $2^k$ sign patterns on the same $k$ variables.
Every assignment falsifies one of these clauses, so their conjunction is
UNSAT. The other $n-k$ variables may remain unused. If $m>2^k$, repeat
clauses to reach the required count; the formula remains UNSAT.
\end{proof}


For example, when $k=2$, every formula with $m\le3$ is SAT, including
$n=m=k=2$. The first possible UNSAT example has four clauses:
\[
 (x\lor y)\land(x\lor\neg y)\land(\neg x\lor y)
 \land(\neg x\lor\neg y).
\]
The requirement of exactly $k$ distinct variables per clause matters.
If ``$k$-SAT'' permits clauses of \emph{at most} $k$ literals, then
$(x)\land(\neg x)$ is already an UNSAT instance with two clauses.

For fixed $k$ and positive density $m/n\to\alpha$, the condition $m<2^k$
eventually ceases to apply. To study this regime, return to the random model
and let $Z$ be its number of satisfying assignments. Every fixed assignment
satisfies a random clause with probability $1-2^{-k}$, so independence gives
\begin{equation}
 \mathbb E Z=2^n(1-2^{-k})^m,
 \qquad \mathbb P(F_k(n,m)\text{ is SAT})\le \mathbb E Z.
 \label{eq:firstmoment}
\end{equation}
Consequently a fixed density
\begin{equation}
 \alpha>\alpha_{\mathrm{FM}}(k):=
 \frac{-\log2}{\log(1-2^{-k})}
 =2^k\log2-\frac{\log2}{2}+O(2^{-k})
 \label{eq:fmthreshold}
\end{equation}
is UNSAT whp. The finite-$n$ probability bound in
\eqref{eq:firstmoment} is explicit; it is still a distributional statement,
not a refutation of a sampled formula. A large expectation does not prove
satisfiability: correlations between assignments can make the second moment
large and concentrate solutions in an atypical minority of formulas.

\subsection{Four different threshold questions}
Write $p_{n,k}(m)=\mathbb P(F_k(n,m)\text{ is SAT})$, a decreasing function
of $m$. A \emph{sharp threshold} means that its change from near one to near
zero occurs over $o(n)$ clauses. A \emph{limiting threshold} additionally
identifies a constant $\alpha_k$ such that
\begin{equation}
 \lim_{n\to\infty}p_{n,k}(\lfloor\alpha n\rfloor)=
 \begin{cases}1,&\alpha<\alpha_k,\\0,&\alpha>\alpha_k.\end{cases}
 \label{eq:threshold}
\end{equation}
Knowing that this constant exists, calculating its value, and bounding the
finite-size transition window are three further, distinct achievements.
Equation~\eqref{eq:threshold} makes no assertion at $\alpha=\alpha_k$.

For $k=2$, the threshold is exactly $1$. The critical scale is
$m=n+\lambda n^{2/3}$, so the clause window has order $n^{2/3}$ and the
density window order $n^{-1/3}$; see
\citet[Theorem 1.1]{bbckw2001}. Implication graphs also give a linear-time
algorithm, including construction of a witness or detection of a contradiction.
Thus a sharp probabilistic transition need not imply computational hardness.

For fixed $k\ge3$, the main milestones are as follows.
\begin{itemize}
\item \citet[Theorem 1.3]{friedgut1999} proved sharpness about a
size-dependent location. That theorem did not prove convergence of the location.
\item \citet[Theorem 1]{dss2022} identified the limiting threshold by the
cavity prediction for all sufficiently large fixed $k$. The unspecified
large-$k$ condition cannot be replaced by $k\ge3$.
\item \citet[Theorem 1.1]{cop2016} proved the large-$k$ asymptotic formula
\begin{equation}
 \alpha_k=2^k\log2-\frac{1+\log2}{2}+o_k(1).
 \label{eq:largek}
\end{equation}
Here the asymptotics concern $k\to\infty$, after the fixed-$k$ threshold
problem. Substituting $k=3$ into the leading terms is not a theorem about
$\alpha_3$.
\item The supplied October 2026 preprint by
\citet[Theorem 1.1 and Corollary 1.2]{carenini2026} states
\eqref{eq:threshold} for \emph{every} fixed $k\ge3$, via a polynomially
sublinear transition window. This resolves the existence question in the
satisfiability conjecture as stated in that paper; it neither evaluates all
$\alpha_k$ nor gives an efficient search algorithm.
\end{itemize}

\subsection{The supplied recent result: precise scope}
In Carenini's proper-clause model, which is our model, the formulas for
different clause counts are constructed from one shared random sequence.
Fix $n\ge k$ and sample independent clauses $C_1,C_2,\ldots$, each uniformly
distributed over the $2^k\binom nk$ proper $k$-clauses. Define
\begin{equation}
 F_k(n,0)=\mathrm{True},\qquad
 F_k(n,m)=C_1\land\cdots\land C_m\quad(m\ge1).
 \label{eq:clauseprocess}
\end{equation}
Thus $F_k(n,m+1)=F_k(n,m)\land C_{m+1}$: all existing clauses are kept
and one new clause is added. For every fixed $m$, the formula has the usual
random $k$-SAT distribution. The term \emph{coupling} specifies their joint
construction on a common probability space. The clauses $C_i$ are independent;
the formulas at different clause counts are dependent because they share
their earlier clauses. Both $n$ and $k$ stay fixed during this process.

Adding clauses can only remove satisfying assignments. Consequently, once
the formula becomes UNSAT, it remains UNSAT. Define the first such clause
count by
\begin{equation}
 T_n=\min\{m\ge0:F_k(n,m)\text{ is UNSAT}\}.
 \label{eq:hittingtime}
\end{equation}
This is a random switching time: different clause sequences can give
different values. Monotonicity gives, for every integer $m\ge0$,
\begin{equation}
 \{T_n>m\}=\{F_k(n,m)\text{ is SAT}\},\qquad
 p_{n,k}(m)=\mathbb P(T_n>m).
 \label{eq:hittingprobability}
\end{equation}
It is finite almost surely, since \eqref{eq:firstmoment} implies
$\mathbb P(T_n>m)\le2^n(1-2^{-k})^m\to0$ as $m\to\infty$ with $n,k$ fixed.
This is the coupling used immediately before Theorem~1.3 of
\citet[p.~2]{carenini2026}; it connects fluctuations of $T_n$ with the
satisfiability transition.

For fixed $k\ge3$ and $0<\eta<1/2$, set
$q_a(n)=\min\{m:p_{n,k}(m)\le a\}$ and
$W_\eta(n)=q_\eta(n)-q_{1-\eta}(n)$.
The preprint gives
\begin{equation}
 \operatorname{Var}(T_n)=O_k(n^{1+2/k}),\qquad
 W_\eta(n)=O_{k,\eta}(n^{1/2+1/k}).
 \label{eq:carenini}
\end{equation}
Thus the density window is $O_{k,\eta}(n^{-1/2+1/k})$; for $k=3$ these
exponents are $5/6$ in clause count and $-1/6$ in density. The hidden
constants depend on $k$ and $\eta$. This does not give a usable sample-size
prescription for a simulation at, say, $n=10^4$.

The geometric object in this proof is a random covering of the Boolean cube:
each clause excludes a codimension-$k$ subcube. The proof relates a cover-time
variance bound to the remaining lifetime of the last nonempty solution set,
then controls that lifetime using forced variables and shorter-clause
comparisons. A convergence criterion for the threshold location completes
the argument; sharpness alone would not suffice. The statement and this proof
structure were checked against the supplied text, but the full argument has
not received an independent review within this project. The source is ECCC
TR26-229, posted 5 October 2026, and is treated here as a recent preprint.

Its predecessor \citep{careniniWindow} gives an $O(n/\log n)$ window.
The stronger polynomial saving in \eqref{eq:carenini} is material to the
limiting-location conclusion, not merely a change in presentation.

Wilson's published lower bound gives $W_\eta(n)=\Omega_{k,\eta}(\sqrt n)$
for fixed $k\ge3$ \citep[Corollary 4]{wilson2002}. Stronger, very recent
manuscripts in the OpenAI mathematics collection claim
$\operatorname{Var}(T_n)=\Theta_k(n)$, for $k\ge4$ and separately $k=3$,
and hence a central clause window of order $\sqrt n$
\citep[Theorem 1.1 in each]{oaVariance,oaVariance3}.
These are separately identified manuscript claims, not a claim of completed
refereeing or formal verification by this survey. They concern fluctuations
about a finite-size center; a width bound alone does not bound the displacement
of that center from $\alpha_k n$ at the same scale.

\subsection{Numerical locations and their evidence}
\begin{table}[htbp]
\centering\small
\begin{tabularx}{\linewidth}{@{}r r Y@{}}
\toprule
$k$ & Location & Interpretation \\
\midrule
2 & $1$ & Exact threshold; published theorem. \\
3 & $\approx4.26675$ & Cavity/population-dynamics prediction, not an exact theorem. \\
4 & $\approx9.931$ & Cavity prediction. \\
5 & $\approx21.117$ & Cavity prediction. \\
6 & $\approx43.37$ & Cavity prediction. \\
7 & $\approx87.79$ & Cavity prediction. \\
10 & $\approx708.936$ & Cavity prediction. \\
\bottomrule
\end{tabularx}
\caption{Useful density landmarks from \citet[Table 1 and Eq.~(43)]{mmz2006},
except for the exact $2$-SAT value. The displayed digits are numerical
predictions, not rigorous error intervals for small fixed $k$.}
\label{tab:locations}
\end{table}

For $3$-SAT, the longstanding published bounds are
$3.52\le\alpha_3\le4.4898$
\citep[Section 5]{kkl2006}\citep[Theorem 1]{diaz2008}.
The constructive lower-bound analysis first supplies a success probability
bounded away from zero; threshold sharpness yields whp existence below the
bound. It should not be paraphrased as a whp guarantee for that algorithm on
every typical formula. The numerical differential-equation checks underlying
the lower bound were not reproduced for this survey.

A preprint by \citet[Theorem 1.1]{vorobyev2026}, posted 6 October 2026,
states that density $4.268$ is UNSAT whp, improving the upper bound if the
argument is accepted. The author supplies a claimed Lean formalization; we
have not built it or checked its correspondence with the probabilistic model.
Even $\alpha_3\le4.268$ would not prove the cavity prediction, nor provide
an UNSAT certificate for a particular formula near that density.

A further September-dated manuscript claims computability of $\alpha_3$:
a finite machine produces a rational approximation with error $2^{-r}$ for
any requested integer $r\ge1$ \citep[Theorem 1.1]{oaComputing}. It supplies
no efficiency bound and does not solve an input SAT instance. It is recorded
as a recent manuscript claim alongside the sharper-window claims.

\subsection{Translating clause models}
Some cited papers sample literal positions independently, allowing repeated
variables within a clause, or use a Poisson clause count. For fixed $k$ and
$m=\Theta(n)$, the probability that \emph{all} clauses have distinct underlying
variables stays bounded away from zero: a clause is proper with probability
$(n)_k/n^k=1-\binom{k}{2}/n+O_k(n^{-2})$. Conditioning on this event
therefore transfers a whp result to our model, though it is not a total-variation
error tending to zero. Poisson and fixed counts can be sandwiched with a fixed
density margin for monotone satisfiability statements. Sampling proper clauses
with or without replacement differs with probability $O_k(n^{2-k})=o(1)$
when $k\ge3$. Finer window statements require their own transfer argument;
Carenini provides one in Proposition 4.2.
